% =====================================================================
%  Climbing Wall -- System Quick Guide (short form, ~2 pages)
%  Compile:  latexmk -pdf system_quickstart.tex
%  Screenshots live in ./images/system_quickstart/
% =====================================================================
\documentclass[11pt,a4paper]{article}
\usepackage{climbingwall}
\graphicspath{{images/system_quickstart/}}
\setdocname{System Quick Guide}

\begin{document}

\guidetitle{Running the Climbing Wall System}%
  {Quick start --- get from a cold setup to live data in a few minutes.}

The system has three parts: the four \textbf{sensor boards} (plugged into the
laptop over USB), the \textbf{backend} (reads the boards, saves recordings,
computes jump height) and the \textbf{website} (the dashboard you use). Set them
up in this order.

% ---------------------------------------------------------------
\section{Plug in the sensor boards}
The four sensor boards connect to the laptop through a \textbf{USB hub}: each
board's USB cable goes into the hub, and the hub plugs into the laptop with a
single cable. There is nothing to log into and nothing to start on the sensor
side --- the backend (next step) reads the boards directly. Boards may be plugged
in before or after the backend starts; they are picked up automatically within a
second or two.

\begin{infobox}
\textbf{First time on a new computer:} set up the Phidget driver once ---
macOS: \textsf{System Settings $\rightarrow$ Privacy \& Security} $\rightarrow$
\textsf{Allow}; Windows: install the Phidget22 drivers from
\code{phidgets.com}. Without it the backend cannot open the boards.
\end{infobox}

% ---------------------------------------------------------------
\section{Start the backend (laptop)}
In a terminal:
\begin{quote}\ttfamily\small
>> cd backend\\
>> source venv/bin/activate\hfill\textrm{\footnotesize(Windows: \texttt{venv\textbackslash Scripts\textbackslash activate})}\\
>> python app.py
\end{quote}
It prints \code{Running on http://127.0.0.1:5001}. Leave it open. One-step
launchers in \code{backend} also handle the first-time setup:
\code{run\_backend.bat} on Windows, \code{./run\_backend.sh} on macOS/Linux
(use the one matching your computer --- \code{.sh} does not run on Windows).

\begin{infobox}
\textbf{Every time you restart the backend} you must afterwards re-click
\textsf{Connect Device} on the website so the dashboard re-attaches to the fresh
stream. Run only \emph{one} backend at a time --- each board can be opened by
exactly one program (this includes the Phidget Control Panel app).
\end{infobox}

% ---------------------------------------------------------------
\section{Start the website (laptop)}
In a second terminal:
\begin{quote}\ttfamily\small
>> cd climbing\_wall\_website\\
>> npm run dev
\end{quote}
Open the printed link (e.g.\ \code{http://localhost:5173}) in a web browser ---
any modern desktop browser works.

% ---------------------------------------------------------------
\section{Connect and record}
\begin{enumerate}
  \item Click \textsf{Connect Device} (top-left). If a warning says only some of
        the 12 sensor channels are attached, check the USB cables of the sensors
        it names.
  \item If the boards have never been calibrated, calibrate them first --- see the
        \emph{Calibration Quick Guide}.
  \item With nothing on the holds, click \textsf{Zero Sensors} to reset the baseline.
        Do this at the start of every session (it is quick and not saved --- see the
        \emph{Calibration Quick Guide}).
  \item Click \textsf{Start Recording} to begin plotting and logging forces.
  \item Climb / hang / push --- watch the charts update live.
  \item Click \textsf{Stop Recording} to pause, then \textsf{Save} (to the backend)
        or \textsf{Export Data} (a \code{.csv} you can open in Excel).
\end{enumerate}

\begin{warnbox}
\textbf{Start a new recording for each climber.} Use \textsf{Stop} then
\textsf{Start New Recording} between users --- this clears the log and keeps the
page from slowing down or crashing on very long sessions.
\end{warnbox}

% ---------------------------------------------------------------
\section{The tabs}
\begin{itemize}
  \item \textbf{Overall Force} --- total pull on each of the four sensors.
  \item \textbf{Force by Direction} --- the $X$, $Y$, $Z$ components per sensor.
        Use \textsf{View in world frame} in the toolbar for tilt-corrected
        vertical / out-of-wall forces (a display-only angle --- it does not change
        the calibration).
  \item \textbf{Recordings} --- open, compare and export the \textbf{five most
        recent} saved sessions. Older ones stay on disk in the backend's
        \code{saved\_recordings} folder but are not listed.
\end{itemize}
A fourth tab, \textbf{Jump Test (BETA)}, is still experimental and is not covered
by these guides.

\begin{figure}[h]\centering
\framedshot{overall_force}{0.97}
\caption{The dashboard: toolbar on top, the four tabs, and the live
\emph{Overall Force} chart.}
\end{figure}

\end{document}
